\section*{Chapitre \ref{ch:b3:bioinformatics} --- Bio-informatique et analyse des séquences}

\begin{solution}{exo:b3:bioinformatics:1}
Global~: les deux séquences alignées d'un bout à l'autre, chaque résidu
dans une colonne --- pour deux protéines que l'on croit homologues sur
toute leur longueur, par exemple des \omterm{def:b3:genomics:comparative}{orthologues} d'une enzyme de
ménage. Local~: la paire de sous-chaînes de meilleur score, le reste
étant négligé --- pour trouver un domaine commun (un domaine SH2 dans
deux protéines de signalisation par ailleurs sans lien), ou un gène dans
une longue séquence génomique.
\end{solution}

\begin{solution}{exo:b3:bioinformatics:2}
Bords $0,-1,-2,-3$ des deux côtés. Ligne A~: $1, 0, -1$. Ligne G~:
$0, 0, -1$. Ligne C~: $-1, -1, 1$. Optimum $F(3,3) = 1$~: AGC sur AAC
sans brèche (correspondance, mésappariement, correspondance~:
$1 - 1 + 1 = 1$).
\end{solution}

\begin{solution}{exo:b3:bioinformatics:3}
$s(a,a) = \lambda^{-1}\log\bigl(q_{aa}/p_{a}^{2}\bigr)$. Le tryptophane
est rare ($p_{W} \approx 0.013$), de sorte que la chance de voir deux
tryptophanes alignés au hasard, $p_{W}^{2}$, est minuscule, et qu'une
paire de tryptophanes conservée est un signe d'homologie bien plus fort
qu'une paire de leucines conservée ($p_{L}
\approx 0.1$)~; le rapport de log-cote est d'autant plus grand.
\end{solution}

\begin{solution}{exo:b3:bioinformatics:4}
La \omterm{thm:b3:bioinformatics:evalue}{valeur E} est le nombre d'\omterm{def:b3:bioinformatics:alignment}{alignements} de score au moins aussi élevé
que l'on attendrait par hasard dans une recherche de cette requête
contre une base de cette taille. $E = 3$ signifie que trois scores de
ce niveau sont attendus par hasard~: la correspondance n'est pas une
preuve d'homologie (elle peut en être une, mais la recherche ne peut
pas le dire).
\end{solution}

\begin{solution}{exo:b3:bioinformatics:5}
$mn = 400\times 2\times 10^{11} = 8\times 10^{13} = 2^{46.2}$. $E(45) =
2^{1.2} \approx 2.3$~; $E(55) = 2^{-8.8} \approx 2\times 10^{-3}$~;
$E(65) = 2^{-18.8} \approx 2\times 10^{-6}$. $E = 10^{-3}$ demande $S' =
46.2 + 10.0 = 56$ bits. Une requête de \num{40} résidus a un $mn$ dix
fois plus petit, $2^{42.9}$~: $53$ bits suffisent --- mais une requête
courte atteint rarement même cela.
\end{solution}

\begin{solution}{exo:b3:bioinformatics:6}
\omterm{def:b3:bioinformatics:motif}{Contenus en information}~: $2$, $1$, $0$, et $2 - H$ avec $H = -(0.7\log_{2}
0.7 + 3\times 0.1\log_{2} 0.1) = 0.36 + 1.00 = 1.36$, donc $0.64$.
Total $R = 3.64$ bits. Correspondances fortuites~: $9.2\times 10^{6}$
positions sur deux brins $\times 2^{-3.64} \approx 7\times 10^{5}$ ---
le \omterm{def:b3:bioinformatics:motif}{motif} est presque inutile à lui seul.
\end{solution}

\begin{solution}{exo:b3:bioinformatics:7}
Une insertion de plusieurs résidus est un seul événement mutationnel,
de sorte que son coût ne devrait pas croître linéairement avec sa
longueur~: un coût d'ouverture plus un petit coût d'extension modélise
cela. Une pénalité d'ouverture trop forte impose des mésappariements là
où il faudrait une brèche et désaligne tout ce qui suit une vraie
insertion~; une pénalité trop faible sème des brèches partout, apparie
des résidus par hasard et gonfle l'identité.
\end{solution}

\begin{solution}{exo:b3:bioinformatics:8}
Pas forcément~: l'\omterm{def:b3:bioinformatics:alignment}{alignement} couvre un segment de \num{130} résidus, ce
qui est la signature d'un domaine commun plutôt que d'un \omterm{def:b3:genomics:comparative}{orthologue}
aligné sur toute sa longueur. Test~: chercher la protéine de mouche en
retour dans le protéome humain (la requête en est-elle la meilleure
correspondance, sur toute la longueur~?), identifier le domaine par un
HMM de \omterm{def:b3:bioinformatics:msa}{profil}, et construire un arbre de gènes de la famille dans
plusieurs espèces.
\end{solution}

\begin{solution}{exo:b3:bioinformatics:9}
Un \omterm{def:b3:bioinformatics:msa}{profil} enregistre, colonne par colonne, ce que la famille tolère~:
une position toujours hydrophobe sans jamais être le même résidu, un
résidu catalytique invariant, une position qui est toujours une brèche
dans la moitié de la famille. Un \omterm{def:b3:bioinformatics:alignment}{alignement} par paires score chaque
résidu contre un seul autre résidu et ne peut pas savoir qu'une valine
en position 40 vaut «~aussi bien~» que l'isoleucine qui s'y trouve dans
la requête. Le \omterm{def:b3:bioinformatics:msa}{profil} pondère en outre les colonnes conservées, de sorte
qu'une faible similitude concentrée là où la famille est conservée
devient significative.
\end{solution}

\begin{solution}{exo:b3:bioinformatics:10}
Avec un \omterm{prop:b3:bioinformatics:logodds}{score espéré} $\mu > 0$ par colonne, le score cumulé le long de
la diagonale de deux séquences aléatoires est une marche aléatoire à
dérive positive~: après $n$ colonnes il vaut environ $\mu n$, de sorte
que le meilleur \omterm{def:b3:bioinformatics:alignment}{alignement local} est pour l'essentiel la totalité, et
que son score croît comme $\mu n$ et non comme $\log n$. La théorie de
Karlin--Altschul, qui exige une dérive négative pour que les scores
élevés soient de rares excursions, ne s'applique pas et aucun $\lambda$
n'existe. Pour de l'ADN à \qty{60}{\%} de GC la chance d'une
correspondance vaut $2(0.3^{2}) + 2(0.2^{2}) = 0.26$, de sorte que le
\omterm{prop:b3:bioinformatics:logodds}{score espéré} vaut $0.26 - 0.74 = -0.48$~: encore négatif, et la
statistique tient~; mais un système tel que correspondance $+1$,
mésappariement $-0.3$ aurait une espérance de $+0.04$ et rendrait le
génome entier comme un seul \omterm{def:b3:bioinformatics:alignment}{alignement}.
\end{solution}

\begin{solution}{exo:b3:bioinformatics:11}
Amorces et chaînage~: trouver par table de hachage les correspondances
exactes ou presque exactes de $k$-mers entre les deux séquences, garder
celles qui s'alignent sur des diagonales cohérentes, les chaîner, et ne
faire de \omterm{thm:b3:bioinformatics:nw}{programmation dynamique} que dans les intervalles entre amorces
chaînées~; on abandonne les \omterm{def:b3:bioinformatics:alignment}{alignements} des régions sans amorce (les
tronçons très divergents). Bande~: si l'on sait les deux séquences
presque colinéaires, ne calculer que les cases d'une bande de largeur
$w$ autour de la diagonale, au coût $wn$ au lieu de $mn$~; on abandonne
tout \omterm{def:b3:bioinformatics:alignment}{alignement} comportant une insertion plus large que la bande.
\end{solution}

\begin{solution}{exo:b3:bioinformatics:12}
Une séquence codante a une période de trois~: les trois positions du
codon ont des compositions en bases différentes (la troisième est la
plus biaisée) et l'usage des codons est inégal dans chaque espèce. Un
modèle à trois états codants en série, chacun émettant des bases avec la
composition de cette position du codon, attribue à l'ADN codant une
probabilité plus élevée que ne le fait l'état non codant, sur une
fenêtre de quelques dizaines de codons, même sans les stops. Dans un
génome humain les exons sont courts (\qty{150}{bp}), séparés par des
introns de kilobases, de sorte que le signal codant est bref et
interrompu~; le modèle doit en outre reconnaître les sites d'épissage,
signaux faibles, et l'énorme quantité de séquence non codante produit
beaucoup de faux segments codants.
\end{solution}

\begin{solution}{pb:b3:bioinformatics:1}
\textbf{1.} Lignes K, Q, T~; colonnes K, A, Q, T~; bords
$0,-1,-2,-3,-4$ et $0,-1,-2,-3$. Ligne K~: $1, 0, -1, -2$~; ligne Q~:
$0, 0, 1, 0$~; ligne T~: $-1, -1, 0, 2$. Optimum $2$~:
\texttt{K-QT} sur \texttt{KAQT}.
\textbf{2.} Meilleur score local $3$~: \texttt{CAT} contre \texttt{CAT}
(résidus 4--6 de GATCAT avec 2--4 d'ACAT)~; \texttt{ATCAT} contre
\texttt{A-CAT} vaut aussi $4 - 1 = 3$.
\textbf{3.} $300\times 450 = 1.35\times 10^{5}$ mises à jour~; contre
la base $300\times 1.2\times 10^{11} = 3.6\times 10^{13}$.
\textbf{4.} $3.6\times 10^{4}$ s, dix heures par requête~; les amorces
de \omterm{def:b3:bioinformatics:blast}{BLAST} sautent presque tout le tableau et répondent en quelques
secondes.
\textbf{5.} $q_{ab} = p_{a}p_{b}e^{\lambda s}$. Alanine~: $e^{1.39} = 4.0$,
$q_{AA} = 0.074^{2}\times 4.0 = 0.022$, rapport $4$. Tryptophane~:
$e^{3.82} = 45$, $q_{WW} = 0.013^{2}\times 45 = 0.0077$, rapport $45$.
Une paire de tryptophanes alignée est $45$ fois plus fréquente chez les
homologues que par hasard, une paire d'alanines quatre fois seulement~;
les paires d'alanines sont néanmoins plus nombreuses en valeur absolue,
parce que l'alanine est fréquente.
\textbf{6.} \qty{24}{\%} tombe dans la zone crépusculaire, où les
\omterm{def:b3:bioinformatics:alignment}{alignements} aléatoires atteignent \qtyrange{15}{20}{\%}~; la \omterm{thm:b3:bioinformatics:evalue}{valeur E}
de l'\omterm{def:b3:bioinformatics:alignment}{alignement}, des \omterm{def:b3:bioinformatics:motif}{motifs} conservés aux bonnes positions, une
correspondance avec un HMM de \omterm{def:b3:bioinformatics:msa}{profil} d'une famille connue, ou un
repliement partagé trancheraient.
\textbf{7.} $mn = 300\times 1.2\times 10^{11} = 3.6\times 10^{13}$~;
$\log_{2}(mn) = 45.0$.
\textbf{8.} $E(92) = 2^{45 - 92} = 2^{-47} \approx 7\times 10^{-15}$~;
$E(38) = 2^{7} = 128$.
\textbf{9.} $E = 10^{-3}$ à $S' = 45 + 10 = 55$ bits~; $E = 1$ à $45$
bits.
\textbf{10.} Dix fois $mn$~: $E \approx 7\times 10^{-14}$, toujours
écrasant.
\textbf{11.} Avec $E = 128$, la dixième correspondance est ce que
produit le hasard~; \qty{40}{\%} d'identité sur $60$ résidus n'est pas
une preuve. Une recherche par \omterm{def:b3:bioinformatics:msa}{profil} des résidus 200--260 dans la base
de domaines pourrait montrer si ce segment est un domaine connu, avec
une statistique dont la comparaison par paires est dépourvue.
\textbf{12.} Des meilleures correspondances réciproques sur toute la
longueur sont compatibles avec une orthologie un pour un~; elles ne la
prouvent pas --- une duplication dans une lignée après la séparation
donne deux \omterm{def:b3:genomics:comparative}{co-orthologues}, et la perte du véritable \omterm{def:b3:genomics:comparative}{orthologue} peut
laisser un \omterm{def:b3:genomics:comparative}{paralogue} comme meilleure correspondance. Le test est un
arbre de gènes portant sur plusieurs espèces.
\textbf{13.} $R = 2 + 2 + 1.6 + 2 + 0.8 + 1.2 + 0.4 + 0.3 = 10.3$ bits.
\textbf{14.} $3.2\times 10^{9}\times 2^{-10.3} \approx 2.5\times 10^{6}$
correspondances fortuites.
\textbf{15.} $\log_{2}(3.2\times 10^{9}/200) = \log_{2}(1.6\times 10^{7})
\approx 24$ bits.
\textbf{16.} Une cooccurrence à espacement fixé ajoute les bits~: $10.3 + 8 =
18.3$, ce qui fournit $8$ des $13.7$ manquants~; environ $5.7$ bits (un
facteur $50$ sur les correspondances fortuites) doivent venir d'ailleurs
--- accessibilité de la \omterm{def:b3:chromatin-epigenetics:nucleosome}{chromatine}, autres partenaires.
\textbf{17.} $H = -(0.5\log_{2}0.5 + 0.5\log_{2}0.5) = 1$ bit~; $R = 2 -
1 = 1$ bit.
\textbf{18.} Un facteur bactérien doit trouver ses quelques sites dans
un génome de \qty{4.6}{Mb} sans l'aide de la \omterm{def:b3:chromatin-epigenetics:nucleosome}{chromatine}~: il lui faut
quelque $19$ bits, et ses sites sont longs et conservés. Un génome
eucaryote est mille fois plus grand, ce qui demande dix bits de plus, et
pourtant ses facteurs ont des sites courts~; ils obtiennent leur
spécificité par la combinaison et par la restriction de la \omterm{def:b3:chromatin-epigenetics:nucleosome}{chromatine}
accessible, ce qui rend aussi la régulation plus évolutive, puisqu'un
site court se gagne et se perd aisément.
\textbf{19.} $3/64 = 0.047$~; le nombre espéré de codons avant un stop
vaut $64/3 \approx 21$.
\textbf{20.} $p_{A} = p_{T} = 0.31$, $p_{G} = p_{C} = 0.19$~: $P(\text{TAA})
= 0.31^{3} = 0.030$, $P(\text{TAG}) = P(\text{TGA}) = 0.31^{2}\times 0.19
= 0.018$~; total $0.066$, longueur espérée du cadre $15$ codons. Un ADN
riche en AT est plein de stops, de sorte que les cadres ouverts
aléatoires sont plus courts et que les longs ressortent davantage.
\textbf{21.} $(61/64)^{100} = e^{-4.80} = 0.008$~; $(61/64)^{300} =
e^{-14.4} = 5.6\times 10^{-7}$.
\textbf{22.} Chaque cadre ouvert maximal se termine sur un stop, et
$3.2\times
10^{9}$ départs de codons contiennent $3.2\times 10^{9}\times 3/64 = 1.5\times
10^{8}$ stops~: environ $1.5\times 10^{8}\times 0.008 = 1.2\times 10^{6}$
cadres aléatoires d'au moins $100$ codons, et $1.5\times 10^{8}\times
5.6\times 10^{-7} \approx 80$ d'au moins $300$.
\textbf{23.} Une bactérie de \qty{4.6}{Mb} compte quelque $4\times 10^{5}$
stops et donc environ \num{3500} cadres fortuits de $100$ codons, mais
presque aucun de $300$~; ses gènes font $300$ codons en moyenne et
\qty{88}{\%} de l'ADN est codant, de sorte qu'un long cadre ouvert est
presque toujours un gène. Chez le ver, \qty{1.5}{\%} de l'ADN code, les
exons font $50$ codons en moyenne --- moins que le seuil fortuit --- et
un million de cadres aléatoires de $100$ codons les submergent. Les
chercheurs de gènes eucaryotes emploient les signaux de sites
d'épissage, le biais d'usage des codons dans un modèle de Markov caché,
l'homologie avec des protéines connues et, surtout, des transcrits
séquencés.
\textbf{24.} Une lecture issue d'un messager épissé s'aligne sur le
génome en deux morceaux séparés par un intron~: la coupure marque les
deux sites d'épissage à la base près~; la couverture en lectures
délimite les exons et les lectures appariées relient les exons
successifs en un seul transcrit, ce qui résout lequel de plusieurs sites
d'épissage candidats est employé.
\textbf{25.} $E \approx 7\times 10^{-15}$ pour la meilleure
correspondance~; environ $24$ bits pour spécifier $200$ sites dans le
génome~; quelque $80$ cadres de lecture fortuits de $300$ codons dans
le génome entier.
\end{solution}
